Проектирование высокопроизводительной криптографической библиотеки требует глубокого понимания математического аппарата, лежащего в основе алгоритмов. В данной главе рассматривается арифметика конечных полей, принципы работы с эллиптическими кривыми, а также архитектурные подходы к аппаратному ускорению вычислений.

\section{Арифметика конечного поля $GF(2^{255}-19)$}

В основе алгоритмов Curve25519 и Ed25519 лежат вычисления в конечном простом поле $\mathbb{F}_p$, где простое число задано как $p = 2^{255} - 19$. Выбор данного числа (псевдо-Мерсенновское простое) Д. Бернштейном \cite{Bernstein2006Curve25519} не случаен: оно оптимизировано для эффективной программной реализации операций редукции по модулю.

\subsection{Выбор представления чисел: Radix-51}

Стандартные 64-битные типы данных недостаточны для хранения 255-битных чисел. Традиционный подход, используемый во многих библиотеках (например, Bouncy Castle), заключается в разбиении числа на 10 элементов (лимбов) по 25.5 бит (Radix-25.5). Однако данный подход требует постоянного приведения типов при работе на 64-битных архитектурах.

В рамках данной работы выбрано избыточное представление \textbf{Radix-51}. Число $X \in \mathbb{F}_p$ представляется в виде полинома:
\begin{equation}
    X = \sum_{i=0}^{4} x_i \cdot 2^{51i}
\end{equation}
где $x_i$ — 64-битные беззнаковые целые числа (\texttt{ulong}).

Каждый элемент хранит 51 бит полезной информации, оставляя 13 бит запаса в 64-битном регистре. Этот запас позволяет выполнять многократные сложения и умножения без риска переполнения и без немедленного выполнения дорогостоящей операции обработки переносов (carry propagation). Данный подход, известный как «отложенный перенос», широко применяется при проектировании аппаратных криптоускорителей (FPGA) \cite{MdpiFPGA} и был адаптирован для программной среды .NET.

Свойство $2^{255} \equiv 19 \pmod p$ позволяет заменять умножение старших разрядов (сдвиги на $2^{255}$) на быстрое умножение на константу 19. Как показывают выкладки, перемножение двух чисел в Radix-51 требует всего 25 скалярных умножений, а возведение в квадрат — 15 умножений, что значительно снижает вычислительную сложность по сравнению с Radix-25.5.

\section{Математика эллиптических кривых и проективные координаты}

Классическая теория эллиптической криптографии строится на кривых в форме Вейерштрасса: $y^2 = x^3 + ax + b$. Алгебраическое сложение точек на такой кривой $R = P + Q$ геометрически интерпретируется как проведение секущей (или касательной в случае $P=Q$) и нахождение третьей точки пересечения \cite{Hankerson2004}.

Формула наклона касательной для удвоения точки имеет вид:
\begin{equation}
    \lambda = \frac{3x_1^2 + a}{2y_1} \pmod p
\end{equation}

Ключевой проблемой данных формул является операция деления. Поиск обратного элемента в конечном поле (например, по Малой теореме Ферма $A^{-1} \equiv A^{p-2} \pmod p$) требует возведения в огромную степень, что эквивалентно сотням умножений.

Для обхода этой проблемы используются \textbf{проективные координаты} $(X, Y, Z)$, где $x = X/Z$, $y = Y/Z$. Использование проективных координат позволяет заменить все операции деления на умножения. Деление выполняется лишь один раз в самом конце алгоритма для перевода результата обратно в аффинные координаты.

\section{Специфика кривых Монтгомери (X25519) и Эдвардса (Ed25519)}

\subsection{Лестница Монтгомери}
Для обмена ключами используется кривая Curve25519, заданная уравнением Монтгомери: $y^2 = x^3 + 486662x^2 + x$. Её главное преимущество — возможность вычислять скалярное произведение $k \cdot P$ с использованием алгоритма «Лестница Монтгомери» (Montgomery Ladder) \cite{RFC7748}. Алгоритм позволяет складывать точки, используя \textbf{только координаты X и Z}, полностью игнорируя координату Y. Это сокращает количество операций и обеспечивает строгое время выполнения (constant-time).

\subsection{Полный закон сложения кривой Эдвардса}
Классические кривые имеют уязвимость: операции сложения равных точек (удвоение) и сложения с нейтральным элементом (бесконечностью) требуют ветвлений (операторов \texttt{if/else}) в коде. Ветвления нарушают принцип \textit{Constant-time} и открывают возможность для атак по времени (Side-Channel Attacks).

В протоколе Ed25519 используется Искривленная кривая Эдвардса (Twisted Edwards Curve): $-x^2 + y^2 = 1 + dx^2y^2$. Её уникальным свойством является наличие \textbf{полного закона сложения} \cite{RFC8032}. Формула сложения двух точек корректно работает для любых точек на кривой, не создавая исключительных ситуаций. Это делает программную реализацию Ed25519 математически невосприимчивой к атакам по времени.

\section{Векторизация и SIMD-ускорение (AVX2)}

В рамках традиционной обработки (в том числе в библиотеке Bouncy Castle) вычисления выполняются последовательно. Поскольку криптографические операции над независимыми ключами не имеют состояния общей памяти (stateless), они идеально подходят для параллельной обработки.

Современные процессоры обладают набором инструкций \textbf{AVX2} (Advanced Vector Extensions), позволяющим оперировать 256-битными векторами. Поскольку в Radix-51 базовым типом является 64-битный \texttt{ulong}, 256-битный регистр процессора вмещает ровно четыре таких элемента.

Это позволяет реализовать концепцию \textbf{Batched Processing} (пакетной обработки). Разработанный алгоритм группирует четыре независимых скалярных умножения (например, 4 процесса генерации ключей) и выполняет их одновременно за счет векторных операций сложения и умножения (\texttt{Vector256<ulong>}). Такой подход позволяет теоретически увеличить пропускную способность системы в 4 раза, что делает разработанное решение оптимальным для высоконагруженных серверных сред.

\section*{Выводы по главе 1}
\addcontentsline{toc}{section}{Выводы по главе 1}

Анализ теоретической базы показал, что для создания высокопроизводительной криптографической библиотеки недостаточно простой трансляции формул на язык программирования. Требуется синергия математических оптимизаций (Radix-51, проективные координаты, полные законы сложения) и аппаратных возможностей процессора (AVX2 Batching). Выбор кривых Curve25519/Ed25519 оправдан их изначальной ориентированностью на безопасную (constant-time) и быструю программную реализацию. На основе этих выводов в следующей главе будет спроектирована архитектура библиотеки для платформы .NET.